\documentclass[12pt]{article}

%polskie
\usepackage[utf8]{inputenc}
\usepackage[T1]{fontenc}
\usepackage[polish]{babel}
\usepackage{amsmath}
%\usepackage{graphics}
\usepackage{graphicx}
\usepackage{newtxtext,newtxmath}

\usepackage{subcaption}
\usepackage{float}

\usepackage[margin=2.5cm]{geometry}
\usepackage{indentfirst}

\usepackage{graphicx}
\usepackage{hyperref}
\usepackage{titlesec}

% \textheight=24cm
% \textwidth=17cm
% \topskip=-10mm
% \topmargin=-20mm
% \leftmargin=-1mm
% \oddsidemargin=-1mm      
% \evensidemargin=-1mm
% \renewcommand{\baselinestretch}{1}
% %\renewcommand{\baselinestretch}{1.2}

\hypersetup{
    colorlinks=true,
    linkcolor=black
}

\begin{document}

\begin{titlepage}
\noindent
\begin{minipage}[c]{90mm}{\includegraphics[width=90mm]{pwr-logo-poziome.png}}\end{minipage} 
\hfill 
\begin{minipage}[c]{60mm}{\includegraphics[width=60mm]{wit-pwr-logo-podpis.png}}\end{minipage}

\vspace{6ex}
\begin{center}
\textbf{\large Inżynieria oprogramowania: projekt}\\[10ex]
\textsf{\Large Temat 1: Wymagania i przypadki użycia systemu informatycznego}\\[5ex]
\textbf{\large STRUSINSKA Małgorzata, (nr indeksu)  \\[1ex]
TORA Kacper, 284474}\\
\end{center}
\vspace{80mm}
\hspace*{100mm}
\begin{minipage}[c]{65mm}
\small%\scriptsize 
Prowadzący:\\
prof. dr hab. inż. Jan Magott
\end{minipage}
\vfill 
\centerline{Wrocław 2026}
\end{titlepage}
\setcounter{page}{2}
\newpage
\tableofcontents
\newpage

\section{Charakterystyka systemu i środowiska wdrożenia}

\subsection{Nazwa i cel projektu}
Projekt dotyczy oprogramowania do zarządzania kontami bankowymi zintegrowanego ze stanowiskiem samoobsługowym.. Głównym celem jest umożliwienie klientom bezpiecznej realizacji operacji gotówkowych (wypłata, wpłata) oraz informacyjnych bezpośrednio przy terminalu, z natychmiastową synchronizacją z serwerem centralnym.

\subsection{Role w systemie}
W projektowanym otoczeniu wyodrębniono dwie główne role wchodzące w interakcję z aplikacją:
\begin{itemize}
    \item \textbf{Klient:} Fizyczny użytkownik korzystający z terminala bankomatowego. Posiada wydaną przez bank kartę płatniczą z kodem PIN i korzysta z ekranu urządzenia w celu wykonania wpłat, wypłat, sprawdzenia stanu konta.
    \item \textbf{Centralny System Bankowy:} Zewnętrzny system banku odpowiedzialny za autoryzację zgłoszeń, prowadzenie bazy rachunków, weryfikację poprawności PIN-u oraz księgowanie operacji finansowych.
\end{itemize}

\subsection{Wymagania formalne i założenia biznesowe}
\begin{enumerate}
    \item Transakcja musi zostać w całości potwierdzona i zaksięgowana albo w pełni wycofana w przypadku awarii podajnika lub braku gotówki.
    \item Architektura systemu musi być modułowa, co pozwoli na bezproblemowe rozszerzenie terminala o dodatkowe kanały (np. wypłaty BLIK).
\end{enumerate}

\subsection{Środowisko techniczne i operacyjne}
\begin{itemize}
    \item Terminal bankomatowy wyposażony jest w moduł dyspensera z czterema kasetami banknotów (nominały 20, 50, 100 oraz 200 PLN) oraz moduł akceptora wpłat.
    \item Komunikacja pomiędzy terminalem a serwerem bankowym odbywa się za pośrednictwem szyfrowanego kanału sieciowego.
    \item Aplikacja projektowana jest z myślą o pracy ciągłej 24/7, a zakładany maksymalny czas odpowiedzi centrali na żądanie autoryzacji nie przekracza 2 sekund.
\end{itemize}

\newpage

\section{Specyfikacja wymagań}

\subsection{Wymagania funkcjonalne}
\begin{itemize}
    \item \textbf{F01:} System odczytuje dane włożonej karty i weryfikuje jej ważność oraz status blokady.
    \item \textbf{F02:} System weryfikuje poprawność wprowadzonego przez klienta 4-cyfrowego kodu PIN.
    \item \textbf{F03:} Karta zostaje automatycznie zablokowana po 3 kolejnych błędnych próbach podania PIN-u.
    \item \textbf{F04:} System umożliwia wypłatę zadanej kwoty gotówki po uprzedniej autoryzacji.
    \item \textbf{F05:} Przed wydaniem środków system sprawdza saldo konta oraz dzienny limit kwotowy klienta.
    \item \textbf{F06:} Terminal dobiera odpowiednią kombinację dostępnych nominałów z kaset i wydaje gotówkę.
    \item \textbf{F07:} System pozwala na wydrukowanie papierowego potwierdzenia po zakończeniu operacji.
    \item \textbf{F08:} Klient ma możliwość podglądu aktualnego salda na ekranie (w PLN oraz w wybranej walucie obcej).
    \item \textbf{F09:} System przyjmuje od klienta pakiet banknotów, przelicza je i zwiększa saldo konta.
    \item \textbf{F10:} Klient ma możliwość zmiany kodu PIN do karty z poziomu menu terminala.
\end{itemize}

\subsection{Wymagania niefunkcjonalne}
\begin{itemize}
    \item \textbf{NF01 (Wydajność):} Średni czas obsługi pojedynczego zapytania autoryzacyjnego wynosi poniżej 2 sekund.
    \item \textbf{NF02 (Bezpieczeństwo):} Dane wrażliwe (kod PIN) są natychmiast szyfrowane i nie są zapisywane w logach aplikacji w postaci jawnej.
    \item \textbf{NF03 (Niezawodność):} Bankomat blokuje możliwość wypłaty, jeśli w kasetach brakuje banknotów do wydania żądanej kwoty.
    \item \textbf{NF04 (Architektura):} Kod tworzony jest zgodnie ze wzorcem MVP (Model-View-Presenter), rozdzielając logikę biznesową od interfejsu.
    \item \textbf{NF05 (Technologia):} Implementacja prowadzona jest w obiektowym języku Java.
    \item \textbf{NF06 (Użyteczność):} Ekran terminala wyświetla komunikaty krok po kroku, a interfejs przystosowany jest do obsługi klawiaturą funkcyjną i ekranem dotykowym.
\end{itemize}

\section{Diagram przypadków użycia}

Poniższy diagram modeluje funkcjonalności udostępniane przez terminal bankomatowy.

\begin{figure}[H]
    \centering
    % Plik wyeksportowany z Visual Paradigm:
    \includegraphics[width=1\textwidth]{Diagram_przypadkow_uzycia.jpg}
    \caption{Diagram przypadków użycia terminala bankomatowego}
\end{figure}

\newpage

\section{Opis scenariuszy przypadków użycia}

\subsection*{PU01. Autoryzacja PIN}
\textbf{Cel:} Identyfikacja klienta i autoryzacja dostępu do konta na podstawie karty i PIN-u.\\
\textbf{Warunki początkowe:} Terminal jest w stanie gotowości; klient włożył kartę do czytnika.\\
\textbf{Warunki końcowe:} Klient uzyskuje dostęp do menu głównego lub karta zostaje zablokowana. Spełniono wymagania: \textbf{F01, F02, F03}.\\
\textbf{Przebieg główny:}
\begin{enumerate}
    \item Terminal odczytuje numer karty i prosi o wprowadzenie 4 cyfr kodu PIN.
    \item Klient wpisuje kod PIN na klawiaturze numerycznej i zatwierdza zielonym przyciskiem.
    \item Terminal przesyła zaszyfrowane zapytanie do systemu.
    \item System bankowy potwierdza zgodność PIN-u i poprawny status karty.
    \item Terminal zeruje licznik nieudanych prób logowania i przechodzi do menu wyboru operacji.
\end{enumerate}

\subsection*{PU02. Wypłata gotówki}
\textbf{Cel:} Pobranie zadanej kwoty w gotówce z powiązanego rachunku bankowego.\\
\textbf{Warunki początkowe:} Klient pomyślnie przeszedł autoryzację (\textbf{PU01}); w kasetach znajdują się banknoty.\\
\textbf{Warunki końcowe:} Pieniądze zostały wydane klientowi, a saldo na koncie zostało pomniejszone. Spełniono wymagania: \textbf{F04, F05, F06}.\\
\textbf{Przebieg główny:}
\begin{enumerate}
    \item Klient wybiera opcję wypłaty i wskazuje kwotę na ekranie.
    \item Terminal weryfikuje w systemie bankowym dostępne środki oraz dzienny limit wypłat.
    \item Algorytm dyspensera wyznacza zestaw nominałów potrzebny do wypłacenia podanej sumy.
    \item System bankowy rezerwuje środki i obciąża rachunek klienta.
    \item Dyspenser wysuwa banknoty w kieszeni wydawczej.
    \item Klient odbiera gotówkę.
    \item Opcjonalnie: następuje rozszerzenie o \textbf{PU04. Wydruk potwierdzenia}.
\end{enumerate}

\subsection*{PU03. Sprawdzenie salda}
\textbf{Cel:} Odczyt i wyświetlenie aktualnych środków na rachunku klienta.\\
\textbf{Warunki początkowe:} Klient jest zalogowany po udanej autoryzacji (\textbf{PU01}).\\
\textbf{Warunki końcowe:} Informacja o saldzie pojawia się na wyświetlaczu. Spełniono wymaganie: \textbf{F08}.\\
\textbf{Przebieg główny:}
\begin{enumerate}
    \item Klient wybiera z menu funkcję „Sprawdź saldo”.
    \item Terminal wysyła zapytanie o stan konta do systemu bankowego.
    \item System bankowy zwraca wartość salda dostępnego oraz ewentualnych blokad.
    \item Terminal prezentuje dane na ekranie.
    \item Opcjonalnie: następuje rozszerzenie o \textbf{PU04. Wydruk potwierdzenia}.
\end{enumerate}

\subsection*{PU04. Wydruk potwierdzenia}
\textbf{Cel:} Wydrukowanie papierowego kwitu z podsumowaniem wykonanej operacji.\\
\textbf{Warunki początkowe:} Klient zrealizował operację bazową (\textbf{PU02, PU03} lub \textbf{PU06}); w drukarce znajduje się papier.\\
\textbf{Warunki końcowe:} Potwierdzenie zostaje wydrukowane i wysunięte z drukarki. Spełniono wymaganie: \textbf{F07}.\\
\textbf{Przebieg główny:}
\begin{enumerate}
    \item Po zatwierdzeniu operacji system wyświetla pytanie o chęć odebrania wydruku.
    \item Klient wybiera opcję „Drukuj potwierdzenie”.
    \item Drukarka generuje kwit z datą, godziną, zamaskowanym numerem karty oraz statusem operacji.
    \item Klient odbiera wydruk.
\end{enumerate}

\subsection*{PU05. Zmiana kodu PIN}
\textbf{Cel:} Nadanie nowego 4-cyfrowego kodu zabezpieczającego dla aktualnej karty.\\
\textbf{Warunki początkowe:} Klient jest zalogowany (\textbf{PU01}).\\
\textbf{Warunki końcowe:} Nowy PIN zostaje zapisany w bazie banku. Spełniono wymaganie: \textbf{F10}.\\
\textbf{Przebieg główny:}
\begin{enumerate}
    \item Klient wybiera z menu opcję zmiany kodu PIN.
    \item Terminal wyświetla prośbę o wpisanie nowego 4-cyfrowego kodu.
    \item Klient wprowadza nowy kod.
    \item Terminal prosi o powtórzenie kodu w celu eliminacji pomyłek.
    \item Klient wpisuje kod ponownie.
    \item Terminal sprawdza zgodność obu ciągów cyfr oraz odrzuca zbyt proste kombinacje (np. cztery jednakowe cyfry).
    \item Nowy kod zostaje przesłany do systemu bankowego, a na ekranie pojawia się komunikat o sukcesie.
\end{enumerate}

\subsection*{PU06. Wpłata gotówki}
\textbf{Cel:} Przyjęcie banknotów przez terminal i powiększenie salda rachunku klienta.\\
\textbf{Warunki początkowe:} Klient jest uwierzytelniony (\textbf{PU01}); wpłatomat jest sprawny technicznie.\\
\textbf{Warunki końcowe:} Gotówka zostaje zdeponowana w kasecie, a konto klienta zasilone. Spełniono wymaganie: \textbf{F09}.\\
\textbf{Przebieg główny:}
\begin{enumerate}
    \item Klient wybiera opcję wpłaty gotówki.
    \item Otwiera się kieszeń depozytowa urządzenia.
    \item Klient wkłada przygotowane banknoty.
    \item Moduł wpłatomatu pobiera plik, sprawdza autentyczność i rozpoznaje nominały.
    \item Terminal prezentuje zestawienie przyjętych banknotów i łączną kwotę wpłaty.
    \item Klient klika „Zatwierdź wpłatę”.
    \item Terminal wysyła polecenie księgowania do systemu bankowego, a środki trafiają na konto.
    \item Opcjonalnie: następuje rozszerzenie o \textbf{PU04. Wydruk potwierdzenia}.
\end{enumerate}

\subsection*{PU07. Sprawdzenie salda w walucie obcej}
\textbf{Cel:} Prezentacja środków na koncie w przeliczeniu na wybraną walutę (dziedziczy zachowanie po PU03).\\
\textbf{Warunki początkowe:} Klient jest uwierzytelniony (\textbf{PU01}); system bankowy udostępnia aktualną tabelę kursów.\\
\textbf{Warunki końcowe:} Na ekranie wyświetla się saldo w PLN oraz jego przeliczona wartość. Spełniono wymaganie: \textbf{F08}.\\
\textbf{Przebieg główny:}
\begin{enumerate}
    \item Klient wybiera opcję podglądu salda w innej walucie.
    \item Terminal wyświetla listę dostępnych walut (EUR, USD, GBP, CHF).
    \item Klient wskazuje wybraną walutę z listy.
    \item System pobiera saldo konta w PLN oraz aktualny kurs średni z systemu bankowego.
    \item Terminal dokonuje kalkulacji i wyświetla równowartość kwoty w wybranej walucie.
    \item Opcjonalnie: następuje rozszerzenie o \textbf{PU04. Wydruk potwierdzenia}.
\end{enumerate}

\end{document}
